% !TeX root = ../papers/nc-p3.tex
\section{相補逆行列の曲面の一般点に至る全窓の六次元族}
\label{sec:split-surface-full-fibers}
\begin{theorem}\label{thm:split-surface-full-fibers}
Bell曲面の一般の $P=(1,-uv,u,v)$ を取り、定理\ref{thm:split-bell-locus}の四つの
指定付き補完の一つを $w_2$ とする。各補完で次が成立する。
\begin{enumerate}
\item 指定法線側の全持上げは、せん断を除くと重み $1,2$ の四座標と三座標のアフィン空間になる。
すべての値で標準Koszul左右AS正則である。原点以外の添字保持同型類は
$\mathbb P(1,1,1,1,2,2,2)$ を残る四元群で割ったものである。
四元群の作用は相補逆行列表示で $P$ を保つ四つの局所基底変換から誘導する。
\item $w_2$ は多重安定で、この六次元族と原点の一類が、
分解も不変直線も指定しない全窓の中で $[w_2]$ に特殊化する同型類を尽くす。
各補完で分解の交換は自己同型により実現する。従って一般の $P$ の四つの補完は
分解を忘れても、$P$ の座標付け替えによる同一視以外には重ならない。
\item 原点の点グラフの純部分は二本のねじれ直線で、各々に四つの長さ一の埋込み部分がある。
非零の持上げの点グラフでは一次元部分は一本の直線で一般被約である。
従ってこの直線は無印で復元できる。曲面基底をも動かす全族の同型類の像は八次元である。
\end{enumerate}
\end{theorem}
\begin{proof}[相補行列、軸交換と閉軌道]
$X=\left(\begin{smallmatrix}0&1\\1&0\end{smallmatrix}\right)$, $Z=\operatorname{diag}(1,-1)$ とする。
第一の四対応を恒等にした表示で、第二は
\[
                  (hZ,h,hZ,h),\qquad h=1,X,Z,XZ
\]
と取れる。実際、両Schur式への代入で逆行列の移送を満たし、
定理\ref{thm:split-bell-locus}の四つの補完を全て与える。
両分解を交換すると第一の四対応は恒等のまま、第二は
$(h,hZ,h,hZ)$ となる。これは対応を表す行列を両端の交換で転置することによる。
新しい交代テンソルのBell四数は逆数である。その四位置に $(1,Z,Z,1)$ を作用させると
旧 $P$ に射影的に戻る。実際 $(01)(23)$ の四数の入替えと第三・第四係数の符号反転で、
$a b=-c d$ を用いる。第一の四対応を保つよう相補側にも対応する逆転置を作用させると、
第二の四対応は、各 $i$ で $k_{i+2}^{-t} L_i k_i^t$ ($k=(1,Z,Z,1)$) となり、
元の $(hZ,h,hZ,h)$ に戻る。各局所行列の共通倍率は相補二切断の復元で吸収できる。
これで全四補完における交換自己同型と、二つの向きの同一視を得る。

複素数上で四数をBell基底に取り、両方の交代テンソルにビットの全反転と全位置の $Z$ が作用すると、
各二文字空間上のGram行列はスカラーとなる。隣接二文字型ではユニタリ行列の積で同じことが従う。
全体の各一文字Gram行列は重みからブロック対角である。六種類の法線二文字の台が全て非零なので、
定理\ref{thm:split-bell-locus}のノルム最小化で二ブロックのスカラーを各位置で一致させられる。
従って $w_2$ は多重安定である。
\end{proof}
\begin{proof}[一般係数体での全障害]
定理\ref{thm:invariant-line-filtered-lifts}の線型条件の核は
$\mathbb Q(u,v)$ 上で20次元、せん断の像は16次元である。
四法線係数の線型写像の階数は13で、その像に、
せん断の余空間の四基底の二次補正の全十係数が既に入る。
\path{computations/verify_split_surface_lifts.py} が四つの $h$ でこの等式を有理関数体上で検算する。
具体的に二つの切断 $M$ で $L,R$ を左右核、$J$ を六次ピボット上の一般逆とし、
三法線係数の基底候補 $U$ を $LUR=0$ で求める。その核にせん断の十六列を入れ、
余空間の四列 $U_b$ を取る。四法線十六列の制約行列に
$L U_b J U_b R$ と $L(U_b J U_c+U_c J U_b)R$ ($b<c$) の両切断の十列を追加しても、
階数は13である。ピボットの選択だけに有理標本を使い、零の判定は一般係数で行う。
これらは有限個の行列の恒等式と非零小行列なので、共通の非空開集合の全点で成立する。

せん断は一次部分で自由に作用するので四座標の切片をちょうど一度通る。
切片上では全ての二次障害が消え、四法線部の一つの二次の持上げに三次元の核を足せる。
残るブロック対角群は共通の法線倍率と、$h$ の移送を保つ $P$ の四元群である。
これは定理\ref{thm:split-inverse-transporter}の安定群計算による。
この群は簡約的なので一次余空間および四法線の線型像の切断を同変に選べる。
共通法線倍率は二群の座標に重み1,2で作用し、表示した商を与える。
正方向のせん断安定群が零であること、軸交換で負方向も零であることから、
射影安定群の恒等成分は上の一次元トーラスだけである。
\end{proof}
\begin{proof}[印を仮定しない被覆]
$w_2$ に特殊化する隣接階数6以下の窓は半安定である。多重安定軌道への
Hilbert--Mumford判定により、一つの基底変更の後に一径数群でこの極限を取れる。
極限を保つ一径数群はその安定群の恒等成分に入るので、
原点自身を除けば上の重みまたはその逆の正の倍である。前者は全持上げで尽くし、
後者は軸交換で前者に戻る。

点グラフについては命題\ref{prop:split-simple-four-roots}の二つの鉛筆を使う。
$u=2,v=3$ では四補完とも各行列式は単純四根、四つの核二次式との共通因子は定数である。
上の検算でこれらの開条件を両座標直線上で確認する。軸以外に曲線はなく、
これらの軸はともに一般被約である。任意の持上げは分岐を持たず、
その点曲線の重複込み次数は2以下となる。実際、非零ファイバーの純曲線を一径数で動かし
射影Hilbertスキームで平坦極限を取れば、極限は基点の点グラフ内に含まれ、
次数は保たれる。二重の指定直線は、その台が一定なので一般被約な同じ直線に二重には特殊化できない。
さらに別の曲線があれば一本の直線であり、指定直線と相補な方の軸に特殊化する。
二直線が交わる条件は閉で、極限の両軸が交わらないので最初から二直線は交わらない。
第二の直線は他に成分がないため四移動不変である。この直線を消す平面を新しい補空間に取ると、
二つの直線のそれぞれについて反対辺の制限が残る二文字平面の行列式関係に入る。
従って法線文字は二個以上かつ二個以下で、テンソルは重み2だけとなる。
補空間の変更はせん断なので、これは切片の原点である。
これで非零点の指定直線の一意性が分かり、放物群による比較が無印でも必要十分となる。
基底の $P$ は閉軌道から有限通りに復元できるので、次元は $2+6=8$ である。
\end{proof}
